\documentclass[11pt]{article}
\usepackage[margin=1in]{geometry}
\usepackage[T1]{fontenc}
\usepackage{lmodern,booktabs,amsmath}
\usepackage[hidelinks]{hyperref}
\usepackage{xurl}
\setlength{\parindent}{0pt}
\setlength{\parskip}{6pt}
\title{Native CARLA Perception-to-Estimator Smoke Test}
\author{Pan Dynamics Collision Avoidance Research}
\date{October 2, 2026}
\begin{document}
\maketitle

\section{Result}
The complete tested input path ran successfully:
native RGB/LiDAR/GNSS/IMU callbacks, causal sensor preparation, actual YOLO26x
car detection, LiDAR association and rectangle fitting, then the unchanged
NRMM estimator. Both the baseline and a camera-blackout run produced 300
finite estimator publications with aligned clocks. However, the target
position and velocity accuracy gates failed. This is a completed coverage
experiment with a failed target-accuracy result, not an end-to-end accuracy
pass. No gain or threshold was changed after truth scoring.

Unlike the preceding state-proxy smoke test, this experiment did not supply
actor position, velocity or acceleration to prediction. Truth was written to
a separate scoring stream. Four prediction-file hashes were frozen before
the scorer opened that stream. The perception summaries contain zero truth
comparisons; the estimator entry point never opens the truth file.

\section{Capture and sensor preparation}
An owned CARLA 0.10.0 instance on port 2011 used off-screen rendering,
low quality, disabled ray tracing and a nice level of 10. The shared port 2000
was not modified, ticked, or signaled. The temporary server was stopped with
SIGINT after a separate query verified zero vehicles and sensors; its launcher
returned the expected signal exit code 130. The original shared process was
absent at the later final check; its subsequent lifecycle was outside this test.

The map was \path{Carla/Maps/Town10HD_Opt}. A map-only search chose spawn
index 105, with 176 m of approximately straight continuation. The physical
Lincoln MKZ ego and Dodge Charger target used the existing waypoint follower,
11 m/s nominal speeds, 10 m look-ahead and a 15 m initial gap. Target speed
convergence differed from the ego: the measured evaluation target speed was
4.757--11.141 m/s, while ego speed was 8.193--11.188 m/s. Their range was
4.729--10.782 m and maximum ego yaw rate was 0.05232 rad/s.

Fifty stationary calibration samples and 150 driving warm-up samples preceded
300 evaluation samples at 0.02 s. All 500 RGB/LiDAR/GNSS/IMU sample groups
matched their CARLA frame IDs, and frame IDs were consecutive. The evaluated
interval was 0--5.98 s. The front camera produced 640-by-640 RGB with a
100-degree field of view and the map postprocess profile. The 64-channel
LiDAR used 1,280,000 points/s, 50 revolutions/s, an 80 m range,
$10/-30$ degree vertical limits, 0.02 m native range-noise standard deviation,
and disabled synthetic dropoff. It collected a full sweep at each world step.

GNSS and IMU were mounted at the vehicle API's declared ego center of mass,
approximately $[0.15,0,0.35]$ m in actor coordinates. Native GNSS noise was
configured to approximately 0.01 m standard deviation per axis. Native IMU
acceleration and gyro standard deviations were 0.01 m/s$^2$ and
0.0005 rad/s, with zero configured biases. LiDAR/IMU seed was 20261002;
GNSS seed was 20261003. These Gaussian models are not deterministic error bounds.

The input adapter uses the map's fixed Mercator georeference to convert native
GNSS to local position. Velocity is the current-endpoint derivative of a
quadratic fit over the last 0.4 s of GNSS positions. Forty stationary IMU
samples initialize roll/pitch; subsequent attitude uses gyro integration and
native compass yaw. CARLA's accelerometer includes gravity, which is removed
before rotation into the planar body frame. CARLA compass has no configured
noise model in this implementation; this ideal heading channel is a scope
limitation. Static RGB/LiDAR extrinsics are shifted from the actor origin to
the ego CG. Estimated sensor poses are composed from native measurements and
those extrinsics, never copied from actor truth transforms. The native
IMU, GNSS and Mercator implementation files in the installed CARLA source
were inspected to determine units, gravity, axes and geographic conversion.

The first capture adapter attempt encountered unsupported legacy camera
gamma/motion-blur attributes. The actual blueprint attributes were enumerated,
the adapter was corrected, and the failed attempt retained. A subsequent
preflight detected an owned vehicle left after that failure; it was explicitly
destroyed after ownership verification. The final capture completed in
26.949 s with no cleanup errors. These setup failures are not successful runs.

\section{Causal perception and estimator connection}
The original perception program uses centered temporal windows and can fill
earlier headings from later frames. A new explicit \texttt{--causal} mode
in \path{perception/yolo_lidar_target_position.py} instead fits each current
cluster against frozen past publications. Missing headings can be carried
forward in the inertial frame and rotated into the current ego frame; startup
headings are never filled from the future. The existing offline default is
preserved. This flag addresses time dependence; the separate sensor-prepared
dataset is what removes truth ego poses from this experiment.

YOLO26x used confidence 0.18, IoU 0.45, a 640-pixel inference size, car-only
acceptance, and the front camera. The rectangle used declared approximate
sedan dimensions of 5.0 by 1.9 m, rather than target truth dimensions.
The smoke configuration used a 1-degree coarse heading step and at most three
current-frame heading iterations; other geometry gates retained their defaults.
No CARLA labels, oracle heading, actor target position, or truth box was used
to choose a detection or fit a position.

The estimator consumed native GNSS position/derived velocity, compensated
native IMU acceleration/gyro and the actual YOLO/LiDAR relative-position
output. Missing perception became unavailable radar-like detections with NaN
positions. Initialization used the first native GNSS/compass sample and first
perception acquisition, with a co-moving, zero-acceleration target hypothesis.
The production configuration, gains and existing RK4 MEX were unchanged.
Every output was checked against the current sample time; the next runtime
was checked at exactly one 0.02-second increment.

For the fault case, only camera images at frames 150--164, corresponding to
$[3.0,3.3)$ seconds, were replaced by a separately generated black test image.
The original capture was preserved. The actual detector and LiDAR association
were rerun; availability was never forced manually at the estimator interface.
All 15 black frames had no target output, and detection resumed at frame 165.
Baseline had one missing localization, frame 269. Blackout had 17 missing
localizations: the 15 injected black frames plus frames 269 and 270.

\section{Accuracy and functional checks}
Accuracy limits were written to \path{protocol.json} before truth scoring.
They retain the preceding smoke test's six estimator RMSE gates. The
post-transient interval contains 200 samples at $t\geq2$ seconds.

\begin{center}
\begin{tabular}{lrrr}
\toprule
Estimator RMSE & Baseline & Blackout & Limit \\
\midrule
Ego position (m) & 0.01379 & 0.01379 & 0.50 \\
Ego inertial velocity (m/s) & 0.04773 & 0.04773 & 0.50 \\
Ego body velocity (m/s) & 0.04699 & 0.04699 & 0.50 \\
Ego yaw (degrees) & 0.05324 & 0.05324 & 5.00 \\
Target relative position (m) & 1.24063 & 1.25972 & 0.50 \\
Target inertial velocity (m/s) & 3.67216 & 3.81549 & 1.50 \\
\bottomrule
\end{tabular}
\end{center}

Baseline valid perception was 299/300 (99.67\%). Blackout valid perception
outside the 15 injected frames was 283/285 (99.30\%). Availability passed
the 95\% gate, but raw perception post-transient position RMSE was
1.24561/1.28333 m and its 95th percentile was 3.19589/3.33940 m, exceeding
the 0.50/1.00 m gates. Maximum raw perception error over the whole capture
was 6.34565 m. Baseline and blackout target-estimator maximum post-transient
errors were 3.72786 and 3.75550 m. Maximum target-position error during the
injected blackout was 0.45243 m; the larger later errors were associated with
other detections, not the blank-camera interval.

All 600 published states were finite. Median estimator-call durations were
approximately 2.30/2.31 ms, with maxima 310.614/310.005 ms. These offline
shared-machine measurements include cold calls; no real-time claim is made.

\section{Failure localization and remaining limits}
The dominant observed failure is object association. The existing
\path{choose_single_detection} sorts cars by confidence and then image area
independently at each frame. The output nevertheless always labels the
measurement with \texttt{targetIdx=1}. Parked roadside cars can therefore
replace the followed moving car without an identity change being reported.

At frame 237, the selected box was approximately
$[370.1,326.5,470.9,382.9]$ pixels, with confidence 0.9471. Inspection shows
that it encloses the parked car to the right, rather than the large red lead
vehicle. Its fitted relative position was $[10.21,2.74]$ m in CARLA body axes,
and its scoring error was 6.135 m. Similar switches appear near frames 9--35
and 243--272. A separate scoring-only projection found 34 valid frames where
the selected box excluded the truth target center. This is a diagnostic,
not a full truth-box-IoU classification. On the 177 post-transient valid
frames whose selected box included that center, raw position RMSE was
0.40927 m. That truth-selected subset is not a corrected deployment result.
Persistent target association remains unimplemented and is the next concrete
issue to resolve; the experiment did not change selection using truth.

The 0.12 m configured target-measurement bound was exceeded on 297 baseline
valid frames and 281 blackout valid frames. Derived GNSS velocity error over
all evaluated samples had RMSE 0.22446 m/s and maximum 1.67199 m/s, exceeding
the configured 0.02 m/s deterministic allowance. Target speed also fell below
the default 10 m/s lower domain. Relative-position bounds were unavailable
at 291/300 frames in both runs. Thus neither the nominal sensor bounds nor a
domain-wide certificate is supported by this native-sensor experiment.

\section{Validation and artifacts}
Five new tests in \path{tests/yoloLidarCausalityTest.py} passed: future-prefix
invariance, no backward filling at stationary startup, ego-frame conversion of
carried headings during dropout, duplicate-time rejection and preservation
of the explicit opt-in/offline default.
The actual baseline and blackout perception outputs matched exactly for all
150 frames before the image change. All estimator fields except runtime
timing also matched exactly over that prefix. A separate native-input probe
changed all subsequent GNSS/IMU samples; the first 150 prepared frames and
poses stayed exactly unchanged. The probe contained an invalid truth file and
a read trap, confirming that preparation never read truth. All clocks, counts,
finite-state checks and independent scoring checks completed successfully.

After the two main runs, a summary-only correction separated unresolved
startup headings from headings actually filled by carrying a past value.
An actual three-frame detector/geometry rerun verified zero filled and three
unresolved headings, with unchanged predictions. The original main-run
summaries are preserved, and source provenance distinguishes the main-run
code hash from the final metadata-only revision.

The execution directory is
\path{simulation_output/carla_full_perception_20261002/}. It retains raw
native sensor data, separate scoring truth, images/clouds, local capture and
adapter scripts, two perception outputs, two estimator outputs and failed
setup attempts. The tracked report directory
\path{report/CARLA_FULL_PERCEPTION_20261002/} records exact commands,
configuration, results, error traces, association/isolation audits and hashes.
Raw capture files, model weights, CARLA tooling and generated binaries are
excluded from the project commit. Identified export copies are retained in
the external archive at
\path{05_Project_B_Collision_Avoidance/2026-10-02_Native_Perception_Estimator/}.
This experiment covers the estimator's front-camera/LiDAR target and native
ego-measurement path; road-curb experiments and closed-loop controller
execution are separate work.
\end{document}
